\documentclass[11pt]{article}
\usepackage{manual_docs_style}
\externaldocument{build/candidate_generation}[candidate_generation.pdf]
\externaldocument{build/cheap_filter_metrics_stage}[cheap_filter_metrics_stage.pdf]
\externaldocument{build/surrogate_confidence_scoring_stage}[surrogate_confidence_scoring_stage.pdf]
\externaldocument{build/generate_question_stage}[generate_question_stage.pdf]

\begin{document}

\docTitle{Dataset Manifest Schema}
\phantomsection\label{docstart:dataset_manifest_schema}

\section*{Overview}
The dataset manifest defines the selected subset of valid, materialized simulator runs that will appear in a TraceBench question bundle (either as test samples or as training samples).
The output of this stage is consumed by the question-generation stage, which turns the selected runs and split metadata into the final TraceBench question bundle.
Selection of the samples is done by two steps.
First, since \name{} reuses the same test samples across matched experimental conditions, selected samples must have confidence above the threshold for each required noise level and matched condition.
Second, the train/test split is done in such a way that accuracy with random metric-based methods is low.

\section*{Invocation}

The selected dataset manifest is produced with:

\begin{DocCode}
env/bin/python workflows/generate/select_dataset.py \
  --model SIMULATOR \
  --policy PRODUCTION_POLICY_ID \
  --dataset DATASET_ID \
  --surrogate-id SURROGATE_ID \
  [--row-slug ROW_SLUG [ROW_SLUG ...]]
\end{DocCode}

The command reads the resolved production run graph, cheap-filter metrics, surrogate-confidence scores, threshold configuration, and selection policy.
It writes the frozen dataset artifacts under:

\begin{DocCode}
models/simulink/<simulator_id>/datasets/<dataset_id>/
\end{DocCode}

including:

\begin{DocCode}
selected_runs.jsonl
sample_manifest.json
selection_report.json
\end{DocCode}

Use \code{--row-slug} when selecting samples for a specific experiment row or set of matched rows.
When multiple rows share the same test set, the selector must enforce the cross-condition rule: each selected test sample must pass the required cheap-filter and surrogate-confidence checks for every noise level and condition in which it will be reused.

\section*{Location}

Selected dataset artifacts live under:

\begin{DocCode}
models/simulink/<simulator_id>/datasets/<dataset_id>/
\end{DocCode}

Recommended files:

\begin{itemize}
\item \code{selected_runs.jsonl}: flat list of selected run IDs and their split/class metadata.
\item \code{sample_manifest.json}: train/test grouping consumed by question generation.
\item \code{selection_report.json}: summary of selection constraints, balance, and rejected candidates.
\end{itemize}

\section*{Inputs}

The selection stage consumes:

\begin{itemize}
\item \code{plans/<production_policy_id>/run_nodes.jsonl}
\item \code{plans/<production_policy_id>/run_edges.jsonl}
\item \code{metrics/<production_policy_id>/cheap_filter_metrics.jsonl}
\item \code{metrics/<production_policy_id>/surrogate_scores/<surrogate_id>.jsonl}
\item \code{surrogates/<surrogate_id>/thresholds.json}
\item \code{experiment_config.json}
\item \code{description_levels.json}
\item \code{selection_policies/<selection_id>.json}
\end{itemize}

\section*{\texorpdfstring{\titlecode{selected_runs.jsonl}}{selected\_runs.jsonl}}

Each line records one selected run.

\begin{DocCode}
{
  "run_id": "run_7fd9daa0055fb3f412ef5e0ab8d4c5ba",
  "family_id": "fam_de4d7c0a5335e9005510cc2442d0971e",
  "split": "test",
  "class_internal": "gravity",
  "class_agent_facing_name": "gravity",
  "production_policy_id": "production_v1",
  "bootstrap_policy_id": "bootstrap_v1",
  "dataset_id": "paper_main_v2",
  "selection_seed": 123,
  "source": "programmatic",
  "validation_profile": "paper_selected",
  "cheap_filter_pass": true,
  "surrogate_filter_pass": true,
  "surrogate_id": "surrogate_ball_drop_v1",
  "noise_level": "high",
  "true_label_confidence": 0.981,
  "confidence_margin": 0.312
}
\end{DocCode}

Allowed \code{split} values are \code{train}, \code{test}, and \code{other}.

\section*{\texorpdfstring{\titlecode{sample_manifest.json}}{sample\_manifest.json}}

\code{sample_manifest.json} groups selected runs into concrete train/test panels for question generation.
It preserves the existing question-generation concept of selecting shared test sets by \code{test_set_slug} and \code{seed}, and train samples by \code{shot_slug}.

Each manifest entry should have the following shape:

\begin{DocCode}
{
  "train_samples": [],
  "train_samples_baselines": [],
  "test_samples": [],
  "test_samples_baselines": [],
  "other_samples": [],
  "seed": 123,
  "test_set_slug": "test_set_main_v1",
  "shot_slug_recipe": "shots_0",
  "train_test_sample_hash": "..."
}
\end{DocCode}

\section*{Selection Rules}

Selection should enforce:

\begin{itemize}
\item only cheap-filter-passing and surrogate-filter-passing intervention runs may be selected as intervention-class samples;
\item only cheap-filter-passing and surrogate-filter-passing baseline runs may be selected as no-intervention samples;
\item no selected production \code{run_id} may appear in the surrogate training manifest for the same \code{surrogate_id};
\item no selected production \code{family_id} may appear in the surrogate training, validation, or calibration split for the same \code{surrogate_id};
\item class counts in each test panel should be as balanced as possible;
\item train and test samples families for the same question must not overlap;
\end{itemize}

\section*{\texorpdfstring{\titlecode{selection_report.json}}{selection\_report.json}}

\code{selection_report.json} must include:

\begin{itemize}
\item timestamp, \termSimulator identifier, production policy ID, dataset ID, selected \code{surrogate_id}, and threshold version;
\item number of production candidates;
\item number passing cheap filter;
\item number passing surrogate confidence filter;
\item number rejected by cheap filter, grouped by failure reason;
\item number rejected by surrogate filter, grouped by low confidence, wrong top-1 surrogate label, low margin, missing score, or wrong noise level;
\item per-class confidence distributions;
\item final selected per-class counts;
\item final pass/fail status for selection acceptance criteria.
\end{itemize}

\end{document}
